% ---------------------------------------------------------------------------
% Validation / evaluation design.
% \input{sections/05_validation} — expects booktabs, and \citep from natbib.
% Numbers here are population counts (final). Accuracy figures live in §Results.
% ---------------------------------------------------------------------------

\section{Validation Design}
\label{sec:validation}

The corpus carries no event annotations, and none exist for it elsewhere: the
documents are guidance, evaluation and reporting material submitted by
individual organisations, not a wire feed with a registry behind it. Any
accuracy claim therefore has to be built rather than looked up. We build it
from three components: a human judging pass over a representative sample of
extracted events (\S\ref{ssec:event-sample}), a second pass over the geocoding
step that turns location strings into footprints (\S\ref{ssec:polygon-sample}),
and a non-LLM baseline judged on the same documents under the same rubric, so
that the system's accuracy has a reference point rather than only an absolute
value (\S\ref{ssec:baseline}).

Two design commitments run through all three. First, every sample is drawn to
mirror its population, not to surface interesting cases; the pipeline also
implements a risk-ranked sampling mode that over-weights hard extractions, and
we report which mode produced each number, because the two answer different
questions --- an estimate versus a lower bound. Second, all sampling is seeded
(42) and the queues regenerate byte-identically from the released scripts, so
the composition of the evaluation set is auditable independently of the
verdicts recorded in it.

\subsection{Event sample}
\label{ssec:event-sample}

The production run over 10{,}000 documents yields 3{,}572 verifiable events
(\S\ref{sec:corpus}). We queue 175 of them for human judgement. Allocation
across the nine hazard families is proportional to each family's share of the
population under largest-remainder rounding, and the draw within a family is
uniform. Proportional allocation, rather than equal cells per family, is what
allows the pooled figure to be read as a population estimate; it costs
precision in the small families, which we address by reporting them as
directional only.

Because the allocation is stratified on hazard family alone, the dimensions it
is \emph{not} stratified on serve as a balance check. The queue is
54.9\% high-confidence against 54.7\% in the population, with a mean
label-free risk score of 0.591 against 0.592; every family cell matches the
population share within 0.3\,pp. We report these rather than asserting
representativeness.

The risk score deserves a note, since it is computed for every event but used
for nothing in the primary analysis. It is a label-free heuristic --- stated
confidence, multi-event documents, vague or year-only dates, country-level
admin tags, thin evidence, non-empty \texttt{missing\_information} --- and in
the alternative sampling mode it drives selection. Under random sampling it is
retained only as a covariate, which lets us ask after the fact whether the
system's own difficulty signals predict human verdicts.

\subsection{Geocoding sample}
\label{ssec:polygon-sample}

Geocoding is validated separately from extraction, because a correct location
string resolved to the wrong footprint and a hallucinated location string are
different failures with different fixes. The unit here is the (event,
location) mention: 5{,}463 in the run, of which 4{,}435 (81.2\%) resolve to a
geometry.

Sampling is two-stage. Stage one draws 400 rows across three strata: 150
\textsc{high} (resolver confidence $\geq$ 0.7, no fallback step, geometry taken
directly from an OSM polygon, and a recognised administrative tier), 150
\textsc{low} (resolved but failing any of those conditions), and 100
\textsc{unresolved} (a mention that produced no geometry at all, recovered by
anti-join against the extraction file). Stage two caps this at a 60-row judging
queue, reallocated in proportion to the underlying pool sizes
(2{,}607 / 1{,}828 / 1{,}028), giving 29 / 20 / 11. The two-stage design keeps
the oversampled thin stratum available for error analysis while making the
judged queue itself representative of all 5{,}463 mentions.

\subsection{Judging protocol}
\label{ssec:protocol}

Events receive one verdict --- \textsc{correct}, \textsc{partial},
\textsc{wrong} --- in a review interface that displays the full extraction
beside the source document. The distinction is about \emph{completeness of a
usable record}, not about field-by-field string agreement:

\begin{itemize}\itemsep2pt
  \item \textsc{correct}: hazard, date and location are all supported by the
    document, and a downstream user would enter this row as-is.
  \item \textsc{partial}: the event is real and correctly identified, but at
    least one field is imprecise, incomplete or over-broad in a way that
    still leaves the record useful --- a year where the document gives a
    month, one of three affected provinces, a hazard label that is right at
    the family level but not the specific one.
  \item \textsc{wrong}: the event is not in the document, is not an event
    (a scenario, an exercise, a projection), or a field is contradicted by
    the source.
\end{itemize}

Verdicts are given against the document, not against the world. An extraction
that faithfully reproduces an inaccurate source is \textsc{correct} under this
rubric; the pipeline is a reading system, and holding it to the accuracy of the
underlying reporting would measure the corpus instead. Two rates follow:
$\text{accuracy} = \text{correct}/\text{judged}$ and
$\text{useful rate} = (\text{correct}+\text{partial})/\text{judged}$. We report
both throughout, since the gap between them is the quantity a deployment
actually negotiates.

Geocoding rows are judged per field rather than holistically ---
\textsc{location} (is this the place the document means?), \textsc{geometry}
(is the returned footprint that place?), \textsc{admin level} (is it at the
right granularity?), \textsc{date} --- with an OpenStreetMap link centred on
the resolved geometry at a zoom scaled to its area, so a country polygon
returned for a district is visible as such on opening.

\subsection{A non-LLM baseline}
\label{ssec:baseline}

Reporting an accuracy figure for an LLM pipeline with nothing beside it invites
the reader to supply their own reference point. We therefore run a classical
extraction baseline over the same input and judge it under the same protocol.

The baseline is spaCy \texttt{en\_core\_web\_lg} NER combined with a hazard
lexicon compiled from the same PreventionWeb label vocabulary the LLM prompt
specifies, so the two runs emit the same \texttt{event\_hazard} strings and
segment identically by family. Hazard mentions are located by pattern; a
$\pm 1$-sentence window around each is scanned for \textsc{gpe}/\textsc{loc}/
\textsc{fac} entities and for \textsc{date} entities that normalise to an
absolute anchor; a window with at least one of each becomes a candidate event,
and candidates sharing a hazard label within a document are merged into one
row, mirroring the one-row-per-distinct-event contract the LLM is held to.

Three things are deliberately held constant. The baseline reads the same
10{,}000 documents (verified by \texttt{asset\_key} overlap against the
production run, since the source query is time-ordered rather than a pinned
list) and the same 12{,}000-character window. It writes the identical output
schema, so the geocoder, the sampling scripts and the review interface run on
it unchanged. And its locations are resolved by the same Nominatim-backed
resolver against the same cache. What differs between the two runs is the
extraction step and nothing else.

\subsection{Paired comparison}
\label{ssec:paired}

The baseline queue is not an independent sample. For each of the 175 events
already queued from the LLM run, we emit the baseline's extraction for the
\emph{same document}, in the same review order, judged with the same rubric in
the same interface. Where a document yields several baseline events, the
counterpart is the one whose hazard family matches the LLM event, falling back
to the highest-confidence candidate; where the baseline found nothing at all,
the row is queued with empty extraction fields and judged as a miss rather than
dropped, since dropping it would silently condition the comparison on the
baseline having produced something.

Pairing makes the difference testable on the discordant pairs. We report an
exact two-sided McNemar test on both binarisations (\textsc{correct}, and
\textsc{correct}$\,\cup\,$\textsc{partial}); exact rather than $\chi^2$ because
at $n=175$ the discordant cells are small, which is where the asymptotic
approximation is least reliable.

One asymmetry remains: metrics defined on the LLM's own positives cannot
express a recall advantage for the baseline. We therefore add a third arm ---
50 events the baseline extracted from documents where the LLM returned nothing,
sampled from that pool and judged under the same rubric. A high correct rate
there is a recall gap in the deployed system; a low one means the baseline's
extra yield is noise.

We also compute automatic agreement between the two runs at corpus scale
(document-level yield, hazard family, location-set $F_1$, date agreement,
geometry IoU). These are reported as agreement, never as accuracy --- neither
run is ground truth. Their role is to be validated against the human verdicts:
if the baseline agrees with the LLM about as often on events humans rejected as
on events humans confirmed, then agreement is not tracking quality and only the
judged arms can rank the two pipelines.

\subsection{Reporting and precision}
\label{ssec:precision}

Every rate is reported with a 95\% Wilson score interval. At $n=175$ the margin
on the pooled event figure is approximately $\pm 6$\,pp; per-family cells range
from roughly $\pm 10$\,pp in the largest family to $\pm 30$\,pp in the
smallest. Family-level numbers are read as directional and the pooled figure is
the one we quote. Sample sizes are pre-declared: the queues are fixed and
ordered before judging begins, judging proceeds in rank order, and we report
the declared $n$ rather than reporting attrition against a larger intended one.

Judging was performed by a single annotator, so no inter-annotator agreement is
available and the boundary between \textsc{partial} and its neighbours carries
that annotator's reading of the rubric above. This is the most substantial
limitation of the evaluation and we return to it in \S\ref{sec:limitations};
the \textsc{useful rate}, which is invariant to exactly where that boundary
falls, is reported alongside accuracy partly for this reason.
